\documentclass{amsart}
% For dealing with references we use the comment environment
\usepackage{verbatim}
\newenvironment{reference}{\comment}{\endcomment}
%\newenvironment{reference}{}{}
\newenvironment{slogan}{\comment}{\endcomment}
\newenvironment{history}{\comment}{\endcomment}

% For commutative diagrams we use Xy-pic
\usepackage[all]{xy}

% We use 2cell for 2-commutative diagrams.
\xyoption{2cell}
\UseAllTwocells

% We use multicol for the list of chapters between chapters
\usepackage{multicol}

% This is generally recommended for better output
\usepackage{lmodern}
\usepackage[T1]{fontenc}

% For cross-file-references

% Package for hypertext links:
\usepackage{hyperref}

% For any local file, say "hello.tex" you want to link to please

% Theorem environments.
%
\theoremstyle{plain}
\newtheorem{theorem}[subsection]{Theorem}
\newtheorem{proposition}[subsection]{Proposition}
\newtheorem{lemma}[subsection]{Lemma}

\theoremstyle{definition}
\newtheorem{definition}[subsection]{Definition}
\newtheorem{example}[subsection]{Example}
\newtheorem{exercise}[subsection]{Exercise}
\newtheorem{situation}[subsection]{Situation}

\theoremstyle{remark}
\newtheorem{remark}[subsection]{Remark}
\newtheorem{remarks}[subsection]{Remarks}

\numberwithin{equation}{subsection}

% Macros
%
\def\lim{\mathop{\mathrm{lim}}\nolimits}
\def\colim{\mathop{\mathrm{colim}}\nolimits}
\def\Spec{\mathop{\mathrm{Spec}}}
\def\Hom{\mathop{\mathrm{Hom}}\nolimits}
\def\Ext{\mathop{\mathrm{Ext}}\nolimits}
\def\SheafHom{\mathop{\mathcal{H}\!\mathit{om}}\nolimits}
\def\SheafExt{\mathop{\mathcal{E}\!\mathit{xt}}\nolimits}
\def\Sch{\mathit{Sch}}
\def\Mor{\mathop{\mathrm{Mor}}\nolimits}
\def\Ob{\mathop{\mathrm{Ob}}\nolimits}
\def\Sh{\mathop{\mathit{Sh}}\nolimits}
\def\NL{\mathop{N\!L}\nolimits}
\def\CH{\mathop{\mathrm{CH}}\nolimits}
\def\proetale{{pro\text{-}\acute{e}tale}}
\def\etale{{\acute{e}tale}}
\def\QCoh{\mathit{QCoh}}
\def\Ker{\mathop{\mathrm{Ker}}}
\def\Im{\mathop{\mathrm{Im}}}
\def\Coker{\mathop{\mathrm{Coker}}}
\def\Coim{\mathop{\mathrm{Coim}}}

% Boxtimes
%
\DeclareMathSymbol{\boxtimes}{\mathbin}{AMSa}{"02}

%
% Macros for moduli stacks/spaces
%
\def\QCohstack{\mathcal{QC}\!\mathit{oh}}
\def\Cohstack{\mathcal{C}\!\mathit{oh}}
\def\Spacesstack{\mathcal{S}\!\mathit{paces}}
\def\Quotfunctor{\mathrm{Quot}}
\def\Hilbfunctor{\mathrm{Hilb}}
\def\Curvesstack{\mathcal{C}\!\mathit{urves}}
\def\Polarizedstack{\mathcal{P}\!\mathit{olarized}}
\def\Complexesstack{\mathcal{C}\!\mathit{omplexes}}
% \Pic is the operator that assigns to X its picard group, usage \Pic(X)
% \Picardstack_{X/B} denotes the Picard stack of X over B
% \Picardfunctor_{X/B} denotes the Picard functor of X over B
\def\Pic{\mathop{\mathrm{Pic}}\nolimits}
\def\Picardstack{\mathcal{P}\!\mathit{ic}}
\def\Picardfunctor{\mathrm{Pic}}
\def\Deformationcategory{\mathcal{D}\!\mathit{ef}}

\setlength{\emergencystretch}{3em}
\begin{document}
\title[Stacks Project, Tag 02KH]{364 --- Flat base change for quasi-coherent cohomology of arbitrary qcqs schemes}
\maketitle
\noindent Copyright (C) 2005--2025 Johan de Jong. Stacks Project authors. Source: \href{https://stacks.math.columbia.edu/tag/02KH}{Tag 02KH}.

\noindent Editorial difficulty note (not author text). Difficulty H1: The separated Cech calculation is insufficient for arbitrary quasi-separated schemes. The proof must compare the Cech-to-cohomology spectral sequences after base change, use the separated comparison on every overlap, and identify the resulting canonical comparison with the higher-direct-image base-change map..

\noindent Editorial provenance note (not author text). History: excerpt from revision \texttt{a0444\allowbreak 6e57e\allowbreak c1fbc\allowbreak 252a8\allowbreak 71afc\allowbreak ec775\allowbreak 2fb28\allowbreak 07b14}, coherent.tex, lines 947--1050. Reference presentation was adapted; original-author context and core support blocks were retained or added. The mathematical wording is unchanged. Licensed under GNU FDL 1.2 or later; no invariant sections or cover texts. See ../licenses/COPYING. Context blocks reproduce author definitions and supporting results; they are not additional counted problems.

\begin{lemma}[Flat base change]
\label{lemma-flat-base-change-cohomology}
Consider a cartesian diagram of schemes
$$
\xymatrix{
X' \ar[d]_{f'} \ar[r]_{g'} & X \ar[d]^f \\
S' \ar[r]^g & S
}
$$
Let $\mathcal{F}$ be a quasi-coherent $\mathcal{O}_X$-module
with pullback $\mathcal{F}' = (g')^*\mathcal{F}$.
Assume that $g$ is flat and that $f$ is quasi-compact and quasi-separated.
For any $i \geq 0$
\begin{enumerate}
\item the base change map of
Cohomology, Lemma \href{https://stacks.math.columbia.edu/tag/02N7}{lemma base change map flat case (Tag 02N7)}
is an isomorphism
$$
g^*R^if_*\mathcal{F} \longrightarrow R^if'_*\mathcal{F}',
$$
\item if $S = \Spec(A)$ and $S' = \Spec(B)$, then
$H^i(X, \mathcal{F}) \otimes_A B = H^i(X', \mathcal{F}')$.
\end{enumerate}
\end{lemma}

\begin{proof}
Using Cohomology, Lemma \href{https://stacks.math.columbia.edu/tag/02N7}{lemma base change map flat case (Tag 02N7)} in (1)
is allowed since $g'$ is flat by
Morphisms, Lemma \href{https://stacks.math.columbia.edu/tag/01U9}{lemma base change flat (Tag 01U9)}.
Having said this, part (1) follows from part (2). Namely,
part (1) is local on $S'$ and hence we may assume $S$
and $S'$ are affine. In other words, we have $S = \Spec(A)$
and $S' = \Spec(B)$ as in (2).
Then since $R^if_*\mathcal{F}$ is quasi-coherent
(Lemma \href{https://stacks.math.columbia.edu/tag/01XJ}{lemma quasi coherence higher direct images (Tag 01XJ)}),
it is the quasi-coherent $\mathcal{O}_S$-module associated to the
$A$-module $H^0(S, R^if_*\mathcal{F}) = H^i(X, \mathcal{F})$
(equality by
Lemma \href{https://stacks.math.columbia.edu/tag/01XK}{lemma quasi coherence higher direct images application (Tag 01XK)}).
Similarly, $R^if'_*\mathcal{F}'$ is the quasi-coherent
$\mathcal{O}_{S'}$-module associated to the $B$-module
$H^i(X', \mathcal{F}')$. Since pullback by $g$ corresponds
to $- \otimes_A B$ on modules
(Schemes, Lemma \href{https://stacks.math.columbia.edu/tag/01I9}{lemma widetilde pullback (Tag 01I9)})
we see that it suffices to prove (2).

\medskip\noindent
Let $A \to B$ be a flat ring homomorphism.
Let $X$ be a quasi-compact and quasi-separated scheme over $A$.
Let $\mathcal{F}$ be a quasi-coherent $\mathcal{O}_X$-module.
Set $X_B = X \times_{\Spec(A)} \Spec(B)$ and denote
$\mathcal{F}_B$ the pullback of $\mathcal{F}$.
We are trying to show that the map
$$
H^i(X, \mathcal{F}) \otimes_A B \longrightarrow H^i(X_B, \mathcal{F}_B)
$$
(given by the reference in the statement of the lemma)
is an isomorphism.

\medskip\noindent
In case $X$ is separated, choose an affine open covering
$\mathcal{U} : X = U_1 \cup \ldots \cup U_t$ and recall that
$$
\check{H}^p(\mathcal{U}, \mathcal{F}) = H^p(X, \mathcal{F}),
$$
see
Lemma \href{https://stacks.math.columbia.edu/tag/01XD}{lemma cech cohomology quasi coherent (Tag 01XD)}.
If $\mathcal{U}_B : X_B = (U_1)_B \cup \ldots \cup (U_t)_B$ we obtain
by base change, then it is still the case that each $(U_i)_B$ is affine
and that $X_B$ is separated. Thus we obtain
$$
\check{H}^p(\mathcal{U}_B, \mathcal{F}_B) = H^p(X_B, \mathcal{F}_B).
$$
We have the following relation between the {\v C}ech complexes
$$
\check{\mathcal{C}}^\bullet(\mathcal{U}_B, \mathcal{F}_B) =
\check{\mathcal{C}}^\bullet(\mathcal{U}, \mathcal{F}) \otimes_A B
$$
as follows from
Lemma \href{https://stacks.math.columbia.edu/tag/02KG}{lemma affine base change (Tag 02KG)}.
Since $A \to B$ is flat, the same thing remains true on taking cohomology.

\medskip\noindent
In case $X$ is quasi-separated, choose an affine open covering
$\mathcal{U} : X = U_1 \cup \ldots \cup U_t$. We will use the
{\v C}ech-to-cohomology spectral sequence
Cohomology, Lemma \href{https://stacks.math.columbia.edu/tag/01ES}{lemma cech spectral sequence (Tag 01ES)}.
The reader who wishes to avoid this spectral sequence
can use Mayer-Vietoris and induction on $t$ as in the proof of
Lemma \href{https://stacks.math.columbia.edu/tag/01XJ}{lemma quasi coherence higher direct images (Tag 01XJ)}.
The spectral sequence has $E_2$-page
$E_2^{p, q} = \check{H}^p(\mathcal{U}, \underline{H}^q(\mathcal{F}))$
and converges to $H^{p + q}(X, \mathcal{F})$.
Similarly, we have a spectral sequence with $E_2$-page
$E_2^{p, q} = \check{H}^p(\mathcal{U}_B, \underline{H}^q(\mathcal{F}_B))$
which converges to $H^{p + q}(X_B, \mathcal{F}_B)$.
Since the intersections $U_{i_0 \ldots i_p}$ are quasi-compact
and separated, the result of the second paragraph of the proof gives
$\check{H}^p(\mathcal{U}_B, \underline{H}^q(\mathcal{F}_B)) =
\check{H}^p(\mathcal{U}, \underline{H}^q(\mathcal{F})) \otimes_A B$.
Using that $A \to B$ is flat we conclude that
$H^i(X, \mathcal{F}) \otimes_A B \to H^i(X_B, \mathcal{F}_B)$
is an isomorphism for all $i$ and we win.
\end{proof}
\end{document}
